% =====================================================================
%  Chapter 24 — Quantum Devices and the Particle
% =====================================================================

\chapter{Quantum Devices and the Particle}
\label{ch:quantum}

The fourth and fifth senses are the most recent, and they are opposite in
character. One of them uses a device that already exists; the other is the
thing the field has been trying to build for twenty-five years.

\section{Quantum devices}The dark matter field is, in a laboratory, an
oscillating classical field with a frequency set by its mass and a strength
set by its coupling. A superconductor qubit is a device whose state is
sensitive to fields at that frequency at that strength.

The theoretical result is that the ambient dark matter field induces a
macroscopic power output in a quantum device, of order
$\rho_\chi g^2/Q$ for a resonator of quality factor $Q$, with a frequency
signature determined by the field's oscillation frequency
\idx{Quantum device dark matter}. The device does not detect a particle; it
detects the field.

This is a genuinely different kind of measurement and it has three
advantages over the conventional ones. The signal is not a rare event, so
the statistics are entirely different. The relevant background is
electromagnetic rather than nuclear, which means the enormous body of
instrumentation developed for radio-frequency and low-noise work applies.
And the observable is a power in a device, which is what a power technology
would need, so the same apparatus serves the science and the engineering.

The limitations are equally clear. The predicted power is small, and the
requirement is that the device's own noise floor be below it, which is a
materials and fabrication problem rather than a physics one. And the
field's coherence across the laboratory is not a problem --- it is a help,
because a single device sees a coherent field rather than a Poisson stream
--- but the field's spatial coherence across a large array is a limit on
how many devices can be combined.

The programme is early. Its merit is that it is the only one in the book
whose success would simultaneously answer a physics question and produce a
technological capability, and the field's funding structures are poorly
suited to that kind of programme.

\section{The particle}The fifth sense is the one the field was built for, and it
remains unachieved.

Twenty-five years of dedicated detectors have produced a set of exclusion
curves that map the parameter space with a precision nothing else in
physics has matched, and not one event that can be attributed to a dark
matter particle. The 248~keV candidate of Chapter~\ref{ch:event248} is the
only live candidate, and it is one event.

The distinction between this sense and the other four is not one of degree.
The gravitational measurements establish that dark matter exists and how
much there is. The quantum device measurements would establish the field's
coupling. The particle detection would establish what the thing is --- its
mass, its spin, its interaction channels --- and everything else follows
from that.

\section{What the gap between them means}The gap between the fourth sense and the
fifth is the gap the field has to close, and it is worth seeing that the two
are not independent.

A quantum device measurement of the field's coupling is a measurement of
something that the particle detection would also determine. If the field is
the QCD axion, the coupling measured in a haloscope and the mass measured
in the same apparatus would together determine $f_a$, and $f_a$ determines
the production mechanism, and the production mechanism determines the
cosmological history of the dark matter.

That is a concrete route from the fourth sense to the fifth, and it is a
route that the field does not currently fund as one. It is also a route that
would settle a good part of the origins question of Part~\ref{part:origins}
without waiting for the particle.

\begin{keytakeaways}
\item A quantum device can respond to the ambient dark matter field as a
coherent signal rather than as a rare event, which changes the statistics and
the background entirely.
\item The predicted power is small, and the limitation is a fabrication and
noise-floor problem rather than a physics one.
\item The particle sense remains unachieved, and it is the only one that
would determine the mass, spin and couplings --- and therefore everything
else.
\item A quantum measurement of the field's coupling is also a measurement of
the parameter that determines the production mechanism, which is a route
from the fourth sense to the fifth that the field does not fund as one.
\end{keytakeaways}
